\abstract{
	% Claude
	% TODO: rewrite this
	Country-level plastic waste reporting is sparse, irregular, and heterogeneous: many nations record only a handful of observations across decades, and indicators differ in unit, cadence, and definitional scope. This data scarcity poses an analytic challenge: how should one forecast or reconstruct national-level trends when fewer than thirty observations are available, and when reporting frequency varies across countries and indicators? In this study, we present (i) PRISM, a large-language-model-assisted extract-transform-load pipeline that turns heterogeneous public sources into structured records, and AUDA, a schema-enforced analytical database derived from PRISM; and (ii) a comparative forecasting and reconstruction benchmark over country-level plastic waste indicators, including plastic waste generation, total resin consumption, and polymer PET consumption. We evaluate naive, statistical, and kernel- and neural-network-based machine learning methods, with hyperparameter selection guided by a one-standard-error rule applied over grid and random search. Across seven experiments spanning 6 to 695 samples, classical and statistical baselines remain competitive with, and on the smallest datasets frequently outperform, more flexible machine learning models, while multivariate neural-network forecasting yields gains when sufficient cross-country information is available. We use weighted absolute percentage error (WAPE) as the primary metric, and discuss the instability of the coefficient of determination on very small holdout sets. The pipeline and benchmark together provide a reproducible foundation for future modeling and policy analysis on plastic waste data.
}

\keyword{machine learning; hyperparameter tuning; plastic data}
